% Weekly exercise sheet for Lecture 13; also included in the synthesis sheet.
\ifdefined\ExercisesSynthesis
\else
\documentclass[11pt]{article}
\usepackage[utf8]{inputenc}
\usepackage[T1]{fontenc}
\usepackage[english]{babel}
\usepackage{lmodern,microtype}
\usepackage[margin=0.9in]{geometry}
\usepackage{amsmath,amssymb,amsthm,mathtools}
\usepackage{enumitem,hyperref}
\newcommand{\E}{\mathbb E}
\newcommand{\Qbb}{\mathbb Q}
\newcommand{\dd}{\,\mathrm d}
\theoremstyle{definition}
\newtheorem{exercise}{Exercise}[section]
\newlist{exparts}{enumerate}{1}
\setlist[exparts,1]{label=\textbf{(\alph*)},itemsep=0.3em,topsep=0.4em,leftmargin=2em}
\newcommand{\lecturesummary}[1]{\begin{quote}\small\textbf{Focus.} #1\end{quote}}
\title{Stochastic Calculus and Option Pricing\\Lecture 13 Exercises}
\author{}
\date{}
\begin{document}
\maketitle
\vspace{-1.2em}
\setcounter{section}{12}
\fi

\section{Week 13 --- Pricing, Hedging, and Model Risk}
\lecturesummary{This variation uses linear algebra to test what a hedge can
span when the model has more than one independent source of risk.}

\begin{exercise}[Count shocks before claiming a perfect hedge]
At a given instant, a model has \(d\) independent Brownian risk factors.
There are \(m\) traded hedging instruments, each with a vector of exposures
to those factors.
\begin{exparts}
  \item Express the local hedging problem as a linear system.
  \item Give the rank condition for being able to cancel every local risk
  direction.
  \item If the rank is smaller than \(d\) and the target has exposure outside
  the span of the instruments, explain why fitting today's derivative price
  cannot remove that unspanned risk.
\end{exparts}
\end{exercise}


\ifdefined\ExercisesSynthesis\else
\end{document}
\fi
